% Moreish v0.2.8 -- reference-matched layout; box typography unified with body.
% Prose unchanged. Presentation only.
\documentclass[10pt,letterpaper]{article}

\usepackage[margin=1in]{geometry}
\usepackage[T1]{fontenc}
\usepackage[utf8]{inputenc}
\usepackage{lmodern}
\usepackage{needspace}
\usepackage{microtype}
\usepackage{amsmath,amssymb,mathtools}
\usepackage{ragged2e}
\usepackage{varwidth}
\usepackage{changepage}
\usepackage{enumitem}
\usepackage[hidelinks]{hyperref}

\setlength{\parindent}{0pt}
\setlength{\parskip}{0.30em}
\linespread{1.02}
\setlength{\abovedisplayskip}{7.5pt plus 2pt minus 2pt}
\setlength{\belowdisplayskip}{7.5pt plus 2pt minus 2pt}
\setlength{\fboxrule}{0.45pt}
\setlength{\fboxsep}{5.5pt}
\setcounter{secnumdepth}{1}
\emergencystretch=2em
\clubpenalty=10000
\widowpenalty=10000
\displaywidowpenalty=10000
\raggedbottom
\hyphenpenalty=350
\exhyphenpenalty=50
\tolerance=1200

\setlist[itemize]{leftmargin=1.5em,itemsep=0.15em,topsep=0.25em}
\setlist[enumerate]{leftmargin=1.7em,itemsep=0.20em,topsep=0.3em}
\providecommand{\tightlist}{%
  \setlength{\itemsep}{0.15em}\setlength{\parskip}{0pt}}

% Centered paper column with conventionally justified prose.
% The body is inset symmetrically into a narrower, centered measure that
% matches the abstract block, so prose, boxes, and displayed equations all
% share ONE column width. Without this the full-width prose fought the
% narrower centered boxes/equations and the text measure appeared to
% expand and contract down the page. Prose stays fully justified inside
% the centered measure (the block is centered, not the individual lines).
\newenvironment{paperbody}
  {\begin{adjustwidth}{0.08\textwidth}{0.08\textwidth}\justifying}
  {\par\end{adjustwidth}}

% Reference-style statement boxes: thin black frame, no fill, natural width
% up to a readable maximum. Each box is centered; text inside is left-aligned.
% Box typography stays in the paper's Latin Modern family; emphasis is weight/size only.
\newsavebox{\voiceboxsave}
\newenvironment{voicebox}
  {\Needspace{3.5\baselineskip}\par\vspace{0.22em}\begin{center}\begingroup
   \begin{lrbox}{\voiceboxsave}\begin{varwidth}{0.82\linewidth}
   \small\rmfamily\upshape\RaggedRight}
  {\end{varwidth}\end{lrbox}\fbox{\usebox{\voiceboxsave}}
   \endgroup\end{center}\vspace{0.22em}}

\newsavebox{\dialogueboxsave}
\newenvironment{dialoguebox}
  {\Needspace{5\baselineskip}\par\vspace{0.26em}\begin{center}\begingroup
   \begin{lrbox}{\dialogueboxsave}\begin{varwidth}{0.84\linewidth}
   \small\rmfamily\upshape\RaggedRight}
  {\end{varwidth}\end{lrbox}\fbox{\usebox{\dialogueboxsave}}
   \endgroup\end{center}\vspace{0.26em}}

\newsavebox{\thesisboxsave}
\newenvironment{thesisbox}
  {\Needspace{3.5\baselineskip}\par\vspace{0.32em}\begin{center}\begingroup
   \begin{lrbox}{\thesisboxsave}\begin{varwidth}{0.82\linewidth}
   \normalsize\rmfamily\bfseries\upshape\RaggedRight}
  {\end{varwidth}\end{lrbox}\fbox{\usebox{\thesisboxsave}}
   \endgroup\end{center}\vspace{0.32em}}

\newsavebox{\plainboxsave}
\newenvironment{plainbox}
  {\Needspace{3.5\baselineskip}\par\vspace{0.24em}\begin{center}\begingroup
   \begin{lrbox}{\plainboxsave}\begin{varwidth}{0.82\linewidth}
   \small\rmfamily\upshape\RaggedRight}
  {\end{varwidth}\end{lrbox}\fbox{\usebox{\plainboxsave}}
   \endgroup\end{center}\vspace{0.24em}}

% Keep section headings from being stranded at page bottoms.
\let\oldsection\section
\renewcommand{\section}{\Needspace{6\baselineskip}\oldsection}
\let\oldsubsection\subsection
\renewcommand{\subsection}{\Needspace{5\baselineskip}\oldsubsection}

\hypersetup{
  pdftitle={This Ontology Is Really Moreish},
  pdfauthor={Leah},
  pdfsubject={Underground Super Hans, Pavlov's Veruca, Jesus Without Root Access, and the Greedy Integral of the Present}
}

\begin{document}
\pagestyle{plain}

\vspace*{-0.32in}
\begin{center}
{\LARGE This Ontology Is Really Moreish\par}
\vspace{0.35em}
{\large Underground Super Hans, Pavlov's Veruca, Jesus Without Root Access,\\
and the Greedy Integral of the Present\par}

\vspace{0.45in}
{\normalsize Leah\par}
\vspace{0.18em}
{\normalsize 16 September 2026\par}
\end{center}

\vspace{0.08in}

\begin{center}
\begin{minipage}{0.82\textwidth}
\begin{center}
\textbf{Abstract}
\end{center}
\justifying
A peculiar failure mode appears when several individually defensible
intuitions are allowed to govern outside their proper scope. A finite
agent cannot perfectly predict its own future. Present experience has an
actuality unavailable to merely projected futures. A life whose value
exists only in deferred promises may fail to be worth inhabiting now.
Moral seriousness sometimes requires absorbing costs rather than
externalizing them onto others. Models should remain corrigible;
authorities should remain questionable; suffering should not
automatically be redeemed merely because it was endured.

None of these claims is obviously absurd. The problem begins when a
truth that works well as a \textbf{sensor} is silently promoted into a
\textbf{governor}.

This paper develops that failure through three exaggerated internal
characters. \textbf{Underground Super Hans} is adversarial epistemic
skepticism: the part that notices when a forecast, institution,
life-plan, or respectable script has mistaken its model for reality.
\textbf{Pavlov's Veruca} is the demand that value actually appear in
lived experience rather than being endlessly deferred. \textbf{The
Jesus/WWJD ideal} is moral symmetry and sacrificial seriousness: the
demand that one ask who is paying for one's comfort.

Each sensor detects something real. Each becomes dangerous when given
global authority.

Their interaction produces what I call the \textbf{Greedy Integral
Problem}: trying to maximize the meaningfulness of a life by maximizing
each moment independently while neglecting the dynamics through which
present actions alter future states, constraints, and available options.
A second bug compounds it: candidate truths learned from religion,
philosophy, literature, skepticism, psychology, or lived observation may
be admitted through phenomenological confirmation while suffering
generated by their application is denied equivalent authority to revise
them. Phenomenology can install the rule more easily than it can
uninstall it.

The worldview becomes \textbf{moreish}: it seasons its own consequences.
Constraint becomes rebellion; rebellion becomes agency; downstream
damage creates more constraint; the new constraint makes the rebellious
frame feel still more correct.

The corrective is not exhaustive planning, religious submission,
respectability, self-denial, or the murder of the sensors. It is
\textbf{authority architecture}: preserve skepticism, present reward,
moral concern, and lived telemetry while denying any one of them sole
executive control. Replace predicted-future worship with adaptive
feedback, option preservation, corrigibility, and suspicion toward
gratuitous irreversible commitments.

\begin{thesisbox}
A sensor may reveal a real truth without being entitled to
become the global policy.
\end{thesisbox}

\end{minipage}
\end{center}

\vspace{0.32em}

\begin{paperbody}
\section{The child philosopher accidentally builds a
controller}\label{the-child-philosopher-accidentally-builds-a-controller}

Begin with a sober premise:

\begin{voicebox}
the map is not the territory.
\end{voicebox}

A finite agent occupies a local position inside a world larger than its
available representation. Future states are not presently available as
actuality; they have to be projected. So far, so good.

The trouble begins when an epistemic distinction quietly compiles into a
normative rule:

\begin{voicebox}
future models are uncertain\\
therefore future models deserve limited authority\\
therefore speculative futures should not dominate the actually lived
present.
\end{voicebox}

Still plausible. Taken seriously, it can even sound faintly ascetic.

It is also an excellent philosophy for accidentally ruining your credit
score.

Then Underground Super Hans wanders into the control room with a can in
his hand:

\begin{voicebox}
\textbf{``Mate, if the future's only a model, why are we letting it ruin
Tuesday? Bit dodgy, innit?''}
\end{voicebox}

The resulting worldview is not merely impulsive. It has an epistemology,
a theory of authenticity, a suspicion of authority, a theory of temporal
value, and eventually a moral system. That is precisely why it is hard
to dislodge.

It has reasons.

\section{The enlightened-person
compiler}\label{the-enlightened-person-compiler}

My personal ontology did not come from one tradition. It was assembled
promiscuously: Jesus and the Bible, Buddhist ideas about attachment and
suffering, Dostoevsky, secular skepticism, Hitchens, Wallace, behavioral
psychology, control theory, physics-shaped intuitions, jokes, and direct
phenomenological experience.

I did not need to believe a source was infallible in order to learn from
it. The procedure was roughly:

\[
\begin{aligned}
        &\text{encounter candidate insight}\\
  \to\;&\text{look for it in lived reality}\\
  \to\;&\text{notice recurrence}\\
  \to\;&\text{internalize}.
\end{aligned}
\]

Jesus appeared to embody something about sacrifice and concern for the
vulnerable. I looked for that structure. Buddhist traditions appeared to
have noticed something about attachment and suffering. I looked for that
structure. Dostoevsky appeared to have noticed recursive consciousness
reopening every closure. I looked for that structure. Hitchens supplied
suspicion toward inherited metaphysical authority. I looked for where
that suspicion was justified.

The source proposed a candidate. Reality got the vote.

That method remains valuable. Its bug was \textbf{scope compilation}.

\[
\begin{aligned}
             &\text{I repeatedly encounter this pattern}\\
  \not\Rightarrow\;&\text{the broadest policy built from it should govern my life}.
\end{aligned}
\]

``Attachment contributes to suffering'' does not imply ``minimize
attachment globally.''\\
``Self-sacrifice can be admirable'' does not imply ``make the self the
preferred sink for moral cost.''\\
``Models are corrigible'' does not imply ``future models deserve almost
no authority.''\\
``Institutions contain bullshit'' does not imply ``institutional
constraints have no causal force.''

The insight may survive. The compiled rule may not.

\begin{plainbox}
\textbf{In plain terms:} finding one excellent line of code does not
mean you paste it into every function.
\end{plainbox}

\section{Three good sensors, three terrible
dictators}\label{three-good-sensors-three-terrible-dictators}

\subsection{1 Underground Super Hans: adversarial
epistemics}\label{underground-super-hans-adversarial-epistemics}

Dostoevsky's Underground Man is useful because reflection has ceased
being an instrument and become an environment. Every closure can be
reopened; every explanation can be attacked from another frame.

Super Hans adds a less tortured instinct:

\begin{voicebox}
\textbf{``Hang on, mate. Have we considered the possibility the whole
setup's an absolute state?''}
\end{voicebox}

This produces \textbf{Underground Super Hans}: the adversarial sensor
permanently occupying the gap between a model and the world.

He detects real things: institutional bullshit, conventions disguised as
necessity, forecasts presented with unjustified certainty, vaguely
specified future rewards, and other people's objective functions
smuggled into the word \emph{reality}.

Tell him:

\begin{dialoguebox}
``Responsible people follow this life-plan.''\\[0.35em]
Hans:\\
\textbf{``Responsible according to who? Sounds like someone's objective
function in a little tie.''}
\end{dialoguebox}

Tell him:

\begin{dialoguebox}
``Everyone knows this is what success looks like.''\\[0.35em]
Hans:\\
\textbf{``Everyone? Correlated sample, mate. Whole thing's
contaminated.''}
\end{dialoguebox}

Tell him:

\begin{dialoguebox}
``If you keep doing this, the consequences will accumulate.''\\[0.35em]
Hans:\\
\textbf{``That sounds suspiciously like a model, mate.''}
\end{dialoguebox}

And that is the problem.

He is formally correct.

It \emph{is} a model.

Hans should remain employed permanently as an auditor. The failure
occurs when teenage philosophy promotes him to governor. Then skepticism
stops auditing forecasts and starts governing how much the future is
allowed to matter at all.

Future Me arrives with a spreadsheet:

\begin{dialoguebox}
\textbf{Future Me:} ``If you repeat this twenty times, I inherit the
accumulated state.''\\[0.35em]
\textbf{Hans:} ``Twenty times? You've invented nineteen imaginary
Tuesdays and want me to suffer for them.''\\[0.35em]
\textbf{Future Me:} ``No.~I'm saying the system has memory.''\\[0.35em]
\textbf{Hans:} ``Allegedly.''
\end{dialoguebox}

Funny. Also expensive.

\subsection{2 Pavlov's Veruca: present
reward}\label{pavlovs-veruca-present-reward}

The second sensor has no patience for epistemology.

She wants it now.

The serious principle underneath her is:

\begin{thesisbox}
A life whose value exists only in projection is not presently a
valuable life.
\end{thesisbox}

That deserves respect. Many supposedly responsible lives are organized
around indefinite deferral: rest after the deadline, art after
stability, freedom after promotion, authenticity after everything
practical is finally handled. The queue never empties.

Veruca asks whether the organism is receiving any actual reward or
merely being trained to chase deferred symbols.

Then a rewarded action becomes easier to repeat:

\[
\text{present resonance}
\to
\text{repeat tendency}.
\]

Technically, that learning shape is closer to \textbf{instrumental or
operant reinforcement} than to classical Pavlovian conditioning.
``Pavlov's Veruca'' is a comic composite, not a technical claim about
conditioning taxonomy.

Her voice is simpler:

\begin{dialoguebox}
\textbf{Veruca:} ``Does this make life better now?''\\[0.35em]
\textbf{Future Me:} ``Not immediately, but--''\\[0.35em]
\textbf{Veruca:} ``Then what exactly are you selling me?''
\end{dialoguebox}

Or:

\begin{dialoguebox}
\textbf{Future Me:} ``We should preserve some money.''\\[0.35em]
\textbf{Veruca:} ``For what?''\\[0.35em]
\textbf{Future Me:} ``Uncertain future states may require liquidity.''\\[0.35em]
\textbf{Hans:} ``Listen to that. `Uncertain future states.' Proper
wizard language.''\\[0.35em]
\textbf{Veruca:} ``Buy the thing.''
\end{dialoguebox}

The deeper mistake is not wanting pleasure. It is compiling

\begin{thesisbox}
present resonance \(\Rightarrow\) good policy.
\end{thesisbox}

Veruca is essential as a \textbf{reward sensor}. She is catastrophic as
head of capital allocation.

\subsection{3 Jesus without root access: moral perturbation
control}\label{jesus-without-root-access-moral-perturbation-control}

The third sensor is morally different.

I do not consider myself conventionally Christian. My closest
description is \textbf{spiritual atheist}. I do not subscribe to
organized religion, and if forced to place a conventional epistemic bet,
I bet my money on atheism.

But that does not exhaust my position.

I cannot establish that reality contains nothing deserving the name God,
divinity, transcendence, or some deeper generative intelligence. So the
more faithful statement is:

\begin{thesisbox}
I bet my money on atheism. I bet my soul on there being
something.
\end{thesisbox}

``Deism-adjacent'' is probably closer than conventional theism, though
even that makes the uncertainty sound cleaner than it feels. I have
something like a personal relationship with a God who may or may not
exist. The uncertainty is part of the relationship.

My spirituality is therefore deliberately cherry-picked. Religious
structures stay when they survive contact with lived experience and
reality as I can understand it; others do not. The Bible is part of that
source history. So are secular critics of it.

None of this came from nowhere. I was Christian-coded rather than
doctrinally raised: baptized, sent to Sunday school early, more explicit
Christianity arriving through my grandparents than through any rigidly
enforced household creed, with Hitchens and the other skeptics showing up
later. None of that stopped the reported figure of Jesus from taking up a
very specific room in my moral imagination.

The Gospel portrayal of Jesus, the Golden Rule, ``love your neighbor as
yourself,'' the wider Christian imagery of carrying burdens, and the
figure of someone willing to absorb cost rather than automatically push
it outward all supplied one candidate invariant:

\begin{thesisbox}
moral seriousness sometimes costs the self something.
\end{thesisbox}

That remains important to me.

The useful sensor asks:

\begin{thesisbox}
``Am I protecting myself by making somebody else carry a cost I
could reasonably help carry?''
\end{thesisbox}

The bug appears when WWJD gets root access.

Combine Golden-Rule symmetry, uncertainty about God, a refusal to grant
my own frame privileged moral status, a physics-shaped intuition that
every action perturbs the environment, and admiration for sacrificial
goodness. The controller can compile an extraordinarily powerful
objective:

\begin{thesisbox}
Drive the negative moral perturbations I produce in the
environment toward zero.
\end{thesisbox}

This was not a metaphorical flourish. It was a genuine control
objective.

And it has a fatal defect: exact moral zero is neither observable nor
obviously coherent. Every choice reallocates burdens. Action has costs;
inaction has costs; boundaries impose costs; rescue has opportunity
costs. A finite agent cannot integrate the complete moral future.

So the optimization never naturally terminates.

And the most available actuator is the self.

\begin{dialoguebox}
\textbf{Moral sensor:} ``Someone is paying for this.''\\[0.35em]
\textbf{Controller:} ``Correct.''\\[0.35em]
\textbf{Moral sensor:} ``Could we bear more of it?''\\[0.35em]
\textbf{Controller:} ``Technically, yes.''\\[0.35em]
\textbf{Hans:} ``Also worth saying, mate, half this `self-care' business
sounds suspiciously like a marketing category.''\\[0.35em]
\textbf{Veruca:} ``And martyrdom is at least dramatic.''
\end{dialoguebox}

Congratulations. You have installed a moral denial-of-service attack.

\begin{dialoguebox}
The legitimate sensor says:\\
``Do not behave as though you are the only person whose suffering
matters.''\\[0.35em]
The pathological governor says:\\
``Therefore behave as though your suffering barely counts.''
\end{dialoguebox}

Those are not the same ethic.

\begin{thesisbox}
WWJD is a moral probe, not root access.
\end{thesisbox}

\section{The Greedy Integral
Problem}\label{the-greedy-integral-problem}

The central mathematical structure is elementary, and it is what earns
the problem the word \emph{integral}.

Picture a meaningful life, schematically, as an accumulation of momentary
value:

\[
J=\int_0^T u(t)\,dt.
\]

The tempting move is to maximize the integrand \(u(t)\) at every instant.
If the moments were independent, that would even work. But a life is not
a bag of independent moments; it is a dynamical system, and each moment
is lived from the state the previous ones leave behind.

Let an agent evolve through discrete time:

\[
x_{t+1}=F(x_t,a_t,\eta_t),
\]

where \(x_t\) is current state, \(a_t\) is action, and \(\eta_t\)
collects uncontrolled influences. Let momentary value be \(u(x_t,a_t)\).
Over horizon \(T\), define total value schematically as

\[
J_\pi(x_0)=
\mathbb E_\pi
\left[
\sum_{t=0}^{T-1}u(x_t,a_t)
\right].
\]

A greedy policy chooses whatever looks best now:

\[
a_t^G\in\arg\max_a u(x_t,a).
\]

When actions do not alter future opportunities, this may be fine. Under
state coupling, however, the current action changes the state from which
later decisions are made. Dynamic programming therefore values
continuation:

\[
a_t^*\in\arg\max_a
\left[
u(x_t,a)+
\mathbb E\,V_{t+1}(F(x_t,a,\eta_t))
\right].
\]

The claim is modest:

\begin{thesisbox}
When present actions alter successor states, maximizing present
reward need not maximize total reward.
\end{thesisbox}

Not ``greedy always fails.'' Not ``delayed gratification is always
superior.'' Not ``the future should dominate the present.''

Only:

\begin{voicebox}
\textbf{the moments are coupled.}
\end{voicebox}

Hans discounts uncertain continuation value. Veruca amplifies current
value. Together they produce:

\begin{voicebox}
``Given the world as it exists right now, choose the most livable option
right now.''
\end{voicebox}

But tomorrow's \emph{right now} is partly inherited from today's action.

The controller keeps inheriting its own output.

\begin{plainbox}
\textbf{In plain terms:} you are not ordering dinner from the same menu
every night. Some choices quietly erase items from tomorrow's menu.
\end{plainbox}

\section{The future is epistemically weak and causally
real}\label{the-future-is-epistemically-weak-and-causally-real}

Present experience arrives through contact. Future experience arrives
through representation.

So Future Me has a peculiar status: causally real, phenomenologically
modeled.

\begin{thesisbox}
The future is epistemically weak and causally real.
\end{thesisbox}

That distinction is the hinge.

\begin{dialoguebox}
The future says:\\
``If this policy continues, the state will probably degrade.''\\[0.35em]
Hans, already opening another can:\\
\textbf{``Probably? There it is. Model language. Absolute scenes.''}
\end{dialoguebox}

He is correct about the representation and wrong about what follows from
that fact.

Weather forecasts are uncertain. Rain still exists. Budget forecasts are
uncertain. Debt still compounds. Relationship models are uncertain.
Repeated behavior still changes relationships.

The uncertainty belongs to the representation. The causal coupling
belongs to the world.

This becomes more important when projection is cognitively expensive.
Detailed planning may require counterfactual simulation, uncertainty
tracking, causal modeling, conflict between futures, suppression of
immediate reward, and action now on behalf of something not yet
phenomenologically present.

Reflection has a budget.

Under exhaustion, overload, threat, financial pressure, pain, or intense
affect, that budget shrinks. In ordinary language: reduced executive
margin under load. In my broader theory of consciousness, TLICA, I use
\textbf{slack} for one family of margins permitting self-directed
intervention after current pressures consume capacity, and
\textbf{Mode-B projection} for reflexive modeling of possible
post-action states.

Nothing here depends on those terms. The familiar phenomena are enough:

\begin{voicebox}
\textbf{Future You is represented. Present You is happening.}
\end{voicebox}

Or, in plainer terms:

\begin{voicebox}
tomorrow is buffering in 240p; today is screaming in 4K.
\end{voicebox}

Of course today wins arguments.

That does not mean today should automatically get more votes.

\section{The epistemic ratchet}\label{the-epistemic-ratchet}

The truth-acquisition procedure described earlier was phenomenological
in spirit:

\[
\text{candidate}
+
\text{recurrent lived confirmation}
\to
\text{internalization}.
\]

The bug was asymmetry.

Suppose I internalize ``attachment contributes to suffering'' because I
repeatedly observe it. Later, a policy derived from that idea produces
suffering. Does that suffering prove the original proposition false?

Strictly, no. Pain is not a refutation, and noticing that is a genuinely
good move.

Then comes the catastrophic extra step:

\begin{voicebox}
``Therefore the suffering supplies no reason to modify the policy.''
\end{voicebox}

Now the system has different admission and revision criteria.

Phenomenological resonance can help install a rule. Phenomenological
damage cannot easily uninstall it because pain is not a logical
refutation.

Thus:

\begin{voicebox}
\textbf{Phenomenology can canonize a rule more easily than phenomenology
can de-canonize it.}
\end{voicebox}

The missing distinction is between two questions:

\begin{enumerate}
\def\labelenumi{\arabic{enumi}.}
\tightlist
\item
  Is the proposition true?
\item
  Is this policy an adequate way to act on it?
\end{enumerate}

Pain may be weak evidence against the first and extremely strong
evidence against the second.

\[
\text{pain}\not\Rightarrow\text{the proposition is false},
\]

but also

\[
\text{pain}\not\Rightarrow\text{keep the controller unchanged}.
\]

\begin{plainbox}
\textbf{In plain terms:} if your smoke detector correctly detects smoke
but electrocutes you every time it goes off, the existence of smoke does
not vindicate the detector.
\end{plainbox}

Distress is not automatically a theorem-refuter.

It is still telemetry.

\section{Why the ontology is
moreish}\label{why-the-ontology-is-moreish}

Underground Super Hans is not merely difficult to refute.

He pays immediately.

Constraint becomes agency:

\begin{voicebox}
``I am trapped'' becomes ``I refuse to grant this thing full
jurisdiction over me.''
\end{voicebox}

The external constraint may remain, but passive imposition becomes
active refusal.

He also pays in epistemic elevation:

\begin{voicebox}
\textbf{``At least we can see the game, mate. Look at the state of
it.''}
\end{voicebox}

And in relief: projected obligations lose force. Future Me stops
prosecuting Present Me.

Most importantly, he pays in comedy.

Once suffering becomes ridiculous, it becomes usable: story, bit,
material, art. The world may still be awful, but now it is at least
fucking funny.

The loop becomes

\[
\begin{aligned}
\text{constraint} &\to \text{Hans frame} \to \text{felt liberation} \\
&\to \text{less correction} \to \text{more constraint} \\
&\to \text{stronger evidence for Hans}.
\end{aligned}
\]

The controller helps generate a world in which its own interpretation
becomes increasingly appropriate.

It is not merely self-sealing.

It is \textbf{self-seasoning}.

The consequences improve the flavor.

\begin{dialoguebox}
\textbf{Future Me:} ``Some of this absurdity is downstream of you.''\\[0.35em]
\textbf{Hans:} ``Even better, mate. Interactive absurdity. Premium
feature.''
\end{dialoguebox}

That is why the ontology is moreish.

\section{Cosmic comedy and the danger of retrospective
meaning}\label{cosmic-comedy-and-the-danger-of-retrospective-meaning}

Consequences are separated from causes in time. Consciousness receives
the current state, not automatically the complete provenance graph.

So one wakes into:

\begin{voicebox}
``Why the fuck is everything like this?''
\end{voicebox}

Only later does part of the chain become visible:

\[
a_{t-n}\to a_{t-n+1}\to\cdots\to x_t.
\]

Then:

\begin{voicebox}
``Oh. Past me did some of this.''
\end{voicebox}

Not all of it. Other people, institutions, chance, material conditions,
and structural constraints remain causal participants. But enough may
become visible for tragedy to acquire a delayed punchline.

Two narrations can then share a phenomenological form:

\begin{dialoguebox}
Providential:\\
``God, you bastard. I see what you were doing.''\\[0.35em]
Self-authorship:\\
``Past me, you bastard. I see what you were doing.''
\end{dialoguebox}

The metaphysics differ. The integration structure may not:

\[
\text{opaque suffering}
\to
\text{expanded context}
\to
\text{retrospective integration}
\to
\text{laughter}.
\]

But this creates another danger: \textbf{redemptive over-coherence}.
Recovery can tempt us to rewrite every wound as necessary. Survival can
masquerade as vindication of the path survived. Comedy can become
theology by accident.

So:

\begin{voicebox}
\textbf{The fact that I can make the arc coherent now does not prove the
arc was necessary, optimal, deserved, providential, or good.}
\end{voicebox}

The joke proves that integration occurred.

It does not prove the suffering was required.

\section{The corrective: invariants, feedback, and option
value}\label{the-corrective-invariants-feedback-and-option-value}

The obvious cure is ``plan more.'' That misses the mechanism.

If long-horizon simulation is expensive under load, demanding more
elaborate simulation asks the depleted controller to solve its own
depletion by working harder.

Likewise:

\begin{voicebox}
``Have more willpower'' names the desired output, not an actuator.\\
``Think about the consequences'' names information already losing the
vote.\\
``Stop being impulsive'' is diagnosis-shaped air.
\end{voicebox}

The useful corrective is an \textbf{adaptive feedback policy} rather
than worship of one predicted future. Schematically:

\[
a_t\in\arg\max_a
\left[
R_{\mathrm{present}}(x_t,a)
+
\lambda\Omega_H(F(x_t,a))
-
\mu I(x_t,a)
\right],
\]

where \(R_{\mathrm{present}}\) is current lived value, \(\Omega_H\) is a
horizon-dependent option or viability functional, \(I\) is irreversible
downside, and \(\lambda,\mu\ge 0\).

The point of \(\Omega_H\) is not raw action count. Two states can offer
the same number of immediate actions while one opens into a rich,
recoverable future and the other into a cul-de-sac. What matters is
future maneuverability: reachability, reversibility, resource margin,
and access to valued states.

This differs from conventional deferred gratification.

Deferred gratification says:

\begin{voicebox}
``Suffer now because I promise a particular future reward.''
\end{voicebox}

Option preservation says:

\begin{thesisbox}
``Do not gratuitously destroy Future You's ability to respond to
whichever future actually arrives.''
\end{thesisbox}

Money preserves options. Mobility preserves options. Skills preserve
options. Sleep preserves options. Administrative maintenance preserves
options. Relationships can preserve options. Avoiding unnecessary
irreversible commitments preserves options.

Hans can audit the forecast without defeating the policy:

\begin{dialoguebox}
\textbf{Hans:} ``You don't know what happens next.''\\[0.35em]
\textbf{Controller:} ``Correct. That's why we're not betting the whole
shop on one outcome.''\\[0.35em]
\textbf{Hans:} ``\ldots Fair play.''
\end{dialoguebox}

And the moral sensor should not be smuggled into an unearned scalar. It
remains a constraint or interrogation:

\begin{voicebox}
Who is carrying the cost?\\
Am I externalizing it?\\
Am I bearing a proportionate share, or using myself as the unbounded
residual sink?
\end{voicebox}

The cheap, robust probes become:

\begin{itemize}
\tightlist
\item
  If I repeat this policy thirty times, where does it tend?
\item
  Is this move reversible?
\item
  Am I buying relief by manufacturing obligation?
\item
  Does this preserve or delete future maneuverability?
\item
  How reliable is the forecast relative to the irreversibility of the
  action?
\item
  Is distress telling me the proposition is false, or that the
  controller implementing it is bad?
\item
  Did I validate a scoped truth and silently compile it into a universal
  rule?
\end{itemize}

These questions do not require prophecy.

They require governance.

\section{The council meeting}\label{the-council-meeting}

Suppose there is a job you hate. Leaving today produces enormous
immediate relief. Staying indefinitely is intolerable. Leaving with no
buffer may destroy housing stability.

The council convenes.

\begin{dialoguebox}
\textbf{Veruca:} ``Quit today. I'm not donating another Tuesday to this
place.''\\[0.35em]
\textbf{Hans:} ``She's got a point. Whole career-progression thing's a
pyramid of lanyards.''\\[0.35em]
\textbf{Future model:} ``Leaving with no buffer creates a serious risk
of deleting several options at once.''\\[0.35em]
\textbf{Hans:} ``Serious risk according to what, mate? Show us the
workings. None of this priesthood-of-the-spreadsheet business.''\\[0.35em]
\textbf{Controller:} ``Good. Stress-test the assumptions.''\\[0.35em]
\textbf{Moral sensor:} ``If we leave, who else absorbs the cost?''\\[0.35em]
\textbf{Controller:} ``Also good. Trace the externality.''\\[0.35em]
\textbf{Veruca:} ``Fine. But if the answer is `suffer indefinitely until
conditions are perfect,' I'm staging a coup.''
\end{dialoguebox}

A policy might survive all four pressures:

\begin{voicebox}
preserve enough resources to exit without deleting critical options,
then leave.
\end{voicebox}

Nobody gets everything.

Nobody gets fired.

No single sensor gets root.

That is governance.

\section{Why this points beyond
itself}\label{why-this-points-beyond-itself}

Everything above can be read without accepting a larger theory. But it
naturally raises deeper questions:

Why does present contact outrank represented futures under some
conditions? Why can someone recognize a rule reflectively yet fail to
deploy the correction when needed? Why does distress sometimes revise a
policy and sometimes merely intensify the frame generating it? Why can a
true observation become maladaptive solely through a change in scope?
How does available control margin alter the probability that reflective
correction can intervene? How does repeated attention rewrite the slow
structures through which later events are interpreted?

These are among the questions my broader project, \textbf{TLICA},
attempts to formalize.

Its vocabulary includes \textbf{Mode-B projection} for reflexive
modeling of possible post-action states, \textbf{slack} for one family
of margins permitting self-directed intervention under pressure,
\textbf{imprinting} for slow modification of the structures through
which experience is read, and \textbf{identity-correlation} for how
tightly content is coupled to the experienced self.

Those terms are not premises of this paper. They are one attempted
formal continuation of the phenomena exposed here.

The intended direction of dependence is:

\[
\text{phenomenon first}
\to
\text{ordinary description}
\to
\text{formal question}
\to
\text{theory},
\]

not

\[
\text{TLICA says it}
\to
\text{therefore it happened}.
\]

This paper should stand even if TLICA later turns out to be wrong.

If it succeeds, however, the reader should now feel why a theory of
consciousness might need to distinguish present contact from projected
possibility; truth from policy; sensor from governor; control margin
from endorsed intention; immediate resonance from long-horizon
viability; lived distress from logical refutation; self-model from self;
and uncertainty from causal irrelevance.

The questions come first.

\section{Conclusion: this ontology is really
moreish}\label{conclusion-this-ontology-is-really-moreish}

The deepest error examined here is not hedonism, rebellion, atheism,
religion, self-sacrifice, distrust of authority, or poor planning.

It is the promotion of \textbf{scoped truths into global governors}.

The future cannot be known perfectly. Therefore it should not dominate
the present by pretending to certainty. But uncertainty does not
dissolve causality.

Present reward is real. A life permanently deferred into projection may
genuinely be malformed. But present resonance does not certify that the
repeated policy preserves the state from which later reward must be
received.

Self-sacrifice can be morally serious. But concern for others does not
require deleting the self from the set of beings whose suffering counts.
The instruction to drive negative moral perturbations toward zero
becomes pathological when the optimization has no observable zero and
quietly assigns the self as the universal actuator.

Phenomenological probing can reveal genuine regularities. But noticing a
pattern does not certify the broadest rule one can compile from it.
Suffering does not automatically refute a proposition; suffering can
still be decisive telemetry that the controller applying it requires
revision.

The child philosopher was not wrong to distrust prophecy. She was not
wrong to notice that attachment can produce suffering. She was not wrong
to admire sacrifice. She was not wrong to test inherited wisdom against
lived experience. She was not wrong to notice that respectable people
often smuggle their objective functions into the word ``reality.''

She was not even wrong to laugh.

The bug was in the compiler.

It erased scope.

It promoted sensors to governors.

It let phenomenological resonance install rules more easily than
phenomenological damage could revise them.

And once those rules began changing the state, their consequences became
new input to the very frame that produced them.

The ontology seasons its own consequences.

It rewards inhabiting itself.

It makes its failures funny.

It turns the damage into material.

It is, in the phenomenological sense developed here,

\begin{thesisbox}
really fucking moreish.
\end{thesisbox}

The answer is not another governor.

It is governance in which no single sensor gets root.

\bigskip

\section*{Notes on sources and intellectual
provenance}

This paper is intentionally synthetic. Its central self-application is
autobiographical and phenomenological; the sources below supplied
concepts, candidate invariants, vocabulary, or cultural figures rather
than independently validating the autobiographical interpretation.

\textbf{Control and decision theory.} Richard Bellman, \emph{Dynamic
Programming} (1957), for continuation value and dynamic programming;
standard feedback-control theory for open-loop versus state-responsive
policy; real-options reasoning as an analogy for preserved
maneuverability under uncertainty.

\textbf{Learning theory.} Ivan Pavlov supplies part of the cultural
imagery. The operative ``rewarded action becomes easier to repeat''
mechanism is closer to instrumental/operant reinforcement; B. F. Skinner
and modern reinforcement-learning literature are closer technical
relatives. Richard Sutton and Andrew Barto, \emph{Reinforcement
Learning: An Introduction}, provide a modern computational vocabulary of
state, action, reward, return, and policy.

\textbf{Literature and cultural figures.} Fyodor Dostoevsky, \emph{Notes
from Underground} (1864), for recursive consciousness and refusal.
\emph{Peep Show}'s Super Hans is used as a comic amplifier of
adversarial epistemic freedom, not as a psychological case study. Roald
Dahl's Veruca Salt supplies the exaggerated demand for immediate
satisfaction. ``Pavlov's Veruca'' is deliberately a composite.

\textbf{Religious and spiritual provenance.} The Gospel portrayal of
Jesus supplied the moral image of costly concern for others and the
vulnerable. The Golden Rule appears in Matthew 7:12 and Luke 6:31.
``Love your neighbor as yourself'' preserves an important symmetry:
\emph{as yourself}, not \emph{instead of yourself}. ``Bear one another's
burdens'' is Galatians 6:2 and belongs to Pauline Christian tradition
rather than the sayings of Jesus. Buddhist traditions concerning
attachment and suffering supplied another candidate structure tested
phenomenologically; their inclusion records intellectual provenance, not
adherence.

\textbf{Secular skepticism and literature of attention.} Christopher
Hitchens contributed to the adversarial suspicion that inherited
metaphysical authority receives no exemption from scrutiny. David Foster
Wallace's writing on attention and ``default settings'' supplied another
secular route into the problem of how repeated framing can govern lived
experience. These sources are inputs to the author's compiler, not
authorities whose presence certifies the resulting ontology.

\section*{Epistemic status}

The control-theory skeleton is standard mathematics used schematically.
The mapping from that skeleton onto the author's lived experience is a
\textbf{structural self-interpretation}, not a clinical diagnosis,
population-level psychological result, or experimentally validated
causal model.

The character constructions are explanatory devices. The theological
discussion describes the author's current position and intellectual
provenance. The proposed corrective - sensor demotion, adaptive
feedback, option preservation, corrigibility, and proportional moral
cost - is a conceptual architecture, not a claim of universal optimality
for arbitrary dynamical systems or arbitrary lives.

The paper's empirical debt therefore remains where it belongs: not in
proving Bellman again, but in any future attempt to establish that this
specific structural interpretation generalizes beyond the authorial
specimen.

\end{paperbody}
\end{document}
